\documentclass[UTF8,a4paper,11pt,fontset=fandol]{ctexart}
\usepackage[margin=2.2cm]{geometry}
\usepackage{amsmath,amssymb,booktabs,longtable,array}
\usepackage[colorlinks=true,linkcolor=blue,urlcolor=blue,citecolor=blue]{hyperref}
\usepackage{xurl}
\usepackage{enumitem}
\setlist{nosep}
\setlength{\parskip}{0.35em}
\setlength{\emergencystretch}{3em}
\title{赛题27：AICore多Stream调度\\约束建模、可复现优化与研究报告}
\author{本地项目研究记录}
\date{2026年10月11日}
\begin{document}
\maketitle
\begin{abstract}
本项目面向AIC/AIV双核池、显存生命周期与跨Stream Event同步共同约束下的DAG调度。
在保持单文件C++提交契约的前提下，以确定性模拟器筛选候选，加入精确候选去重、串行快路径、
小图显存动态规划、最后消费者释放优先级和图局部性回退。
主体开发集与不同种子的holdout各143条有可行证据的记录，冻结基线分别通过107和105条，
优化版均通过143条。共同可行案例的完工时间与同完工时间显存峰值没有回退。
随后扩展资源争抢、密边、独立并行oracle、多消费者交叠与递归共享图。
最终1193条案例记录1124按预期通过，69条有独立可行witness但启发式仍未解；
记录包含重复官方样例、参数控制及一条预期不可行输入，不能当作1193个独立随机可行拓扑。
979条严格回归、144份独立最优标尺、单文件编译及四项CTest通过。
新增严格半容量资源互斥界使12条解析控制的线程CPU合计中位数下降82.8\%，
143条通用集的CPU反而上升9.3\%，专项收益与负结果分别报告。
这些结果证明本地流程可运行、且覆盖率有所提高；它们不构成对未知官方数据的全局最优或有解必达保证。
\end{abstract}

\section{赛题契约与研究边界}
官方题面链接见文献\cite{official}，最近核对的更新时间为2026年10月8日16:17。
输入包括硬件核数、Stream数、显存上限、Event记录/等待耗时，以及ID可稀疏、记录可乱序的算子DAG。
节点类型固定为AIC或AIV；核数请求须在1至128内且不超过对应核池，单节点计算耗时在1至$10^9$内。
累计时间采用64位整数。输出接口保持
\texttt{ScheduleCvGraph(const HardwareInfo\&, const vector<OperatorNode>\&, int64\_t)}。
提交源仅包含\texttt{problem.h}与C++标准库；JSON适配器、Python生成器和独立验证器不进入提交文件。
本地\texttt{problem.h}是题面转录，尚未取得官方原文件与测试包；平台时间/内存限制未公开。

每个Stream顺序执行Wait、Compute、Record。全部等待事件就绪后，Wait阶段才开始，耗时为事件数量乘等待单价；
Record耗时独立于计算资源。非负Event编号只有一个记录者，每个Record至少被一次Wait使用，单次等待列表不能有重复，
禁止同Stream Event同步。允许的额外Event只改变竞争时序，不改变原始数据依赖与内存释放。
本地方案对每条跨Stream直接数据边，在消费者自身等待列表中显式同步；这是保守实现选择，
不声称官方判题必然排斥已有同步传递所提供的等效依赖。

令算子$i$的计算开始、结束时刻为$a_i,e_i$，工作区、输出大小为$W_i,O_i$，直接后继集合为$D_i$。
输出释放时刻及内存为
\[
 \ell_i=\begin{cases}\max_{j\in D_i} e_j,&D_i\ne\varnothing,\\T,&D_i=\varnothing,\end{cases}
 \qquad M(t)=\sum_i W_i\mathbf1_{[a_i,e_i)}(t)+\sum_i O_i\mathbf1_{[a_i,\ell_i)}(t).
\]
区间左闭右开，同一时刻先处理完成与释放，再按压缩后的Stream编号尝试启动；某Stream资源不足不阻止后续Stream。
空Stream先删除并压缩编号。目标为$T=\max_i e_i$，不含最后一次Record尾部耗时；
首先要求$T\leq\mathrm{deadline}$，再按$(T,\max_t M(t))$字典序优化。
这使得普通拓扑排序、只按最早完成分流、以及把额外Event误作内存释放依赖，都可能产生错误答案。

\section{经典工作及其可迁移部分}
HEFT采用向上秩排序和最早完成处理器选择，是异构DAG列表调度的重要基线\cite{heft}。
本项目借用关键路径秩与最早完成分流思路；赛题的节点类型与核数请求固定，并具有统一显存和确定的Event成本，
所以本实现并不是完整HEFT复现，也不继承论文实验中的性能结论。

资源受限项目调度RCPSP的串行/并行调度生成方法见Kolisch\cite{kolisch}，
Kolisch与Hartmann的计算研究指出不同生成方法与局部搜索组件需在受控预算下比较\cite{hartmann}。
本项目据此保留多优先级、多Stream数以及邻域交换，但将提交侧精确模拟作为最终合法性门槛。
赛题的输出存活内存不是恒定可再生资源，其状态会随最后一个直接消费者完成而改变，不能直接照搬标准RCPSP资源表。

Marchal等研究通过附加虚拟依赖限制内存并行度，并讨论最大拓扑割、显存约束和关键路径之间的关系\cite{memorydag}。
本项目的额外Event guard受到这一思路启发，但其数据模型、释放规则与赛题不完全相同；
本地所有guard始终和原始数据边分开保存。任务树的有限内存调度研究\cite{memorytree}
还提示树形、汇聚与任意DAG应分开考虑，不能用少数规则拓扑代替工作负载覆盖。

CUDA官方文档说明Stream等待已记录事件所捕获的工作；CUDA允许重复记录覆盖Event状态，未记录Event还可表示空工作\cite{cuda}。
赛题规定唯一记录者与初始未就绪，因此本地模拟必须按题面运行，不能把CUDA运行时语义当作判题规则。
Ascend CANN 6.2.RC2文档说明RecordEvent与StreamWaitEvent用于多Stream同步，实际硬件Event还有资源管理要求\cite{ascend}。
该版本文档仅作为工程机制参考；本项目当前不执行真实NPU任务，也不推断其他CANN版本的限制。

\section{算法及实现批次}
\subsection{候选搜索和精确筛选}
图验证后计算关键路径秩、平均释放潜力和无权高度。基础组合包含6种优先级（超过2048节点时4种）、
Stream数集合中的1、2、3、4、8、16、32、64、128和允许的最大值（本地上限256），以及两种类型亲和性。
候选通过完整Wait/Compute/Record状态机判断资源、显存、死锁和deadline。
512节点以内另有算子迁移与相邻无直接依赖交换。基础唯一评估预算为220、96、36，分别对应
$n\leq512$、$512<n\leq2048$、$n>2048$；后续策略使用独立有界预算，避免挤掉原有搜索结果。
若没有通过验证的候选，接口显式抛出异常，不返回未经验证的计划。

\subsection{降低重复与模拟成本}
第一批完整计划签名虽然释放了重复候选预算，但增加事件映射与排序成本；该负结果已保留。
最终签名在计划生成之前构造：压缩Stream的任务序列决定所有原始数据同步，
另外只记录去重后的真实跨Stream guard边。采用哈希桶加完整向量比较，哈希碰撞不会错误剪枝。
相同Event生产者的编号重命名不改变题面规定的统一等待成本，因而可归一化。
串行候选没有Event与计算重叠，采用线性精确内存扫描；优先级排序只计算一次。
提交生成的稠密Event用向量映射，外部任意稀疏Event保留通用映射。

\subsection{显存可行性和图局部性}
当静态串行策略均不能满足显存且$n\leq18$时，执行子集动态规划。
状态$S$为已完成算子集合，存活输出为
\[
 L(S)=\sum_{i\in S,\ D_i=\varnothing\ \lor\ D_i\not\subseteq S}O_i.
\]
对所有前驱已在$S$中的节点$i$，转移
\[
 P(S\cup\{i\})=\min\max\{P(S),L(S)+W_i+O_i\}.
\]
这精确求解串行最小峰值，复杂度$O(n2^n)$，两个主要状态数组在18节点时约3MiB；
并行最优性仍未求解。另有动态列表优先级仅在某节点成为生产者的最后直接消费者时计入释放收益，
并试验净释放、工作区和高度的不同组合。

重尾时长、窄内存和随机ID的局部DAG使上述策略仍可能失败。
新增回退从无向依赖邻接建立低度数BFS局部性排列，采用3次远端探索，
随后以该排列及逆序作为Kahn拓扑排序优先级。它不使用ID中的生成顺序，也不使用输入附带的witness。
这一局部性回退覆盖开发集和独立种子中的剩余困难图。多不连通分量的BFS采用epoch标记，
不再对每个分量反复清空全图距离数组。
保守资源预留与完成前沿guard构造仍作为后备：提前预留Wait阶段资源会牺牲并行度，
这批实验对已收集困难图没有额外恢复作用，因此放在局部性回退之后。

\subsection{探索中的负结果}
将新内存优先级插到原portfolio之前，曾导致共同案例6个完工时间回退，且CPU成本上升，故未采用该顺序。
动态最后消费者优先级、无权高度及保守资源列表本身，均未解决首轮剩余7个重尾窄内存失败。
依赖局部DFS也观察到峰值超过witness上限。这些失败与各批原始JSON/CSV保留在报告目录，
最终方案采用独立追加预算、先精确去重，再按缺解情况启用局部性回退。

\section{工作负载与独立测试方法}
原压力集153条，包含3个官方示例、四类DAG的规模梯度、零/大Event、乱序稀疏ID、64位累计时间、128核请求、
紧并行deadline、空图和1个可证明无解deadline。
开发与holdout各143条，由不同固定种子20261011、20271011生成；每集含3个官方例、
37个8/12/16/18节点精确串行DP oracle、60个32/128/512/2048节点随机结构解析并行witness、
27个蝶形/梳形/宽汇聚witness、12个1024/4096/10000节点显式串行witness，
以及4个96/768/3072/9216节点的不连通窄内存汇聚图。
随机结构涵盖稀疏任意前驱、局部不规则邻域、fork/join、弱连接链群和按前驱数量加权的随机连接。
其中历史case标签\texttt{scale\_free}仅为本地名称，未证明其度数分布具有幂律性质。
时长含重尾，内存含零值和大工作区，Event含0、3、$10^6$成本。

所有可行证据均由Python构造，不查询被测solver。串行证据按直接消费者计数解析内存；
小图oracle枚举可达子集；并行证据按Stream、数据边与额外guard递推时刻，
硬件核池容量大于等于Stream数乘单任务最大核数，保证其解析时刻不被核争抢改变。
随后由独立C++模拟器验证witness的合法性、deadline、解析$T$及peak，
再和提交侧通用模型对照；返回方案另外重复独立验证。
原压力集提供低容量核争抢场景；解析并行witness集的核容量保证有意避开核争抢，不能混为全面核争抢覆盖。

新增80轮性质检查，对固定计划进行节点ID变换、输入记录逆序、稀疏Event重命名、Wait列表逆序及空Stream插入，
要求合法性、$T$和peak保持一致。密集Event快速路径的返回指标必须等于独立模拟。
模型/parser检查共265项，另有72项语义检查（包含固定种子随机图）。这些提供模型一致性的有限证据，
不是形式证明。每套JSON为本地\texttt{aicore-local-v1}适配格式，绝不冒充官方\texttt{.case}格式。
第十九批额外加入4项独立常量期望的整数归一边界检查，模型/parser/算术共269项通过，未将算术检查计为DAG样本数。

\section{实测结果与成本}
冻结基线源哈希以\texttt{5ef41afab1c08366}开头，下表主体配对版本源以\texttt{6729fe58d8cdb0e0}开头，
最近完成验证的提交源以\texttt{b68148bfed40aa1c}开头。完整哈希及版本范围保存在各paired摘要中。
测量设备为AMD Ryzen 9 8945HX，16物理核/32逻辑线程，31.3GiB RAM，Windows 11 build 26200。
编译环境为GCC 15.2.0、CMake 4.2.0、C++17，Release使用\texttt{-O3 -DNDEBUG}；单文件另以\texttt{-O2}独立编译。
solver自身为单线程CPU工作，构建可并行4任务；本流程不需要GPU/NPU、CANN或真实显存分配。

每集各执行3组交替顺序的before/after测试。
\path{solve_us}用\path{steady_clock}测\path{ScheduleCvGraph}调用墙钟，
包含操作系统抢占与等待，不是线程实际消耗的CPU时间。
下文历史批次表格中的``CPU''均指\emph{CPU求解器的调用墙钟时间}，旧数字不变，不与之后新增线程CPU字段混用。
全CLI墙时另外包含JSON读取、witness验证、模型对照和结果输出。失败的基线案例没有被记作$T=0$参与质量比较，
共同案例定义为前后均能返回合法deadline方案。host后台负载未做CPU绑定，故报告中位数与原始重复数据，
不将微小计时差异解释为稳定的硬件加速保证。

\begin{table}[ht]
\centering\small
\begin{tabular}{lrrrr}
\toprule
指标 & 开发基线 & 开发优化 & holdout基线 & holdout优化\\
\midrule
可行记录/143 & 107 & 143 & 105 & 143\\
共同案例数 & \multicolumn{2}{c}{107} & \multicolumn{2}{c}{105}\\
共同solver总时中位数(s) & 2.0769 & 1.9569 & 2.3777 & 2.3791\\
完整CLI墙时中位数(s) & 4.730 & 4.560 & 4.334 & 4.943\\
共同案例$T$改善数 & -- & 11 & -- & 12\\
同$T$峰值改善数 & -- & 10 & -- & 9\\
$T$/同$T$峰值回退数 & -- & 0/0 & -- & 0/0\\
\bottomrule
\end{tabular}
\caption{主体配对版本对冻结基线，数据来自两个paired-final摘要及其每次原始结果。}
\end{table}

原153条压力记录继续153/153通过；最终三套合计439条case-record验证通过，其中共享官方样例重复存在，
不称为439个独立唯一问题。研究两套共286条均有构造性可行证据，基线仅212条返回可行方案，
新增恢复74条。最终三个CTest均通过，单文件独立编译通过，未发现模型T/peak不一致。
第一个官方示例的手工方案$T=365$、peak=320得到复现；本地solver其余两个官方例采用不同合法分流，
分别得到$T=122$和$150$、peak=55和45，不能与手工方案$T=124$和152混淆。

\begin{table}[ht]
\centering\small
\begin{tabular}{rrr}
\toprule
节点数 & 优化solver中位数(ms) & 观测最大值(ms)\\
\midrule
18 & 0.762 & 0.766\\
128 & 9.058 & 16.530\\
512 & 25.668 & 63.625\\
2048 & 37.475 & 82.158\\
4096 & 48.288 & 127.748\\
9216 & 54.694 & 58.981\\
10000 & 96.965 & 207.998\\
\bottomrule
\end{tabular}
\caption{主体配对版本开发集3轮按节点数汇总；包含失败基线对应的恢复案例。}
\end{table}

研究集配对完整CLI最高工作集约405.4MiB，含约24.7MB输入JSON的通用DOM、原始字符串、
witness、模拟器与结果数据；这不是solver单独内存，也不是题面中的device peak。
三轮开发和三轮holdout合计12次完整CLI墙时约56.2秒，
数据生成需要Python标准库，18节点oracle为指数成本，只适用于小图。
模拟推进到下一个完成/等待/记录时刻，$10^9$量级时长不会按微秒逐点循环；原压力集累计时间达到$4.096\times10^{12}$。
CMake、CTest、JSON/CSV和最终可重放失败输入均本地保存。各批失败文件可能非空，它们代表历史负结果；
第六批主体paired的after失败文件均为空；后续探索性困难集的真实失败分别保留并报告。
截至第十八批，项目目录含1593个文件、约1.39GiB磁盘数据，主要保留各批输入、返回计划、原始重复测量和冻结二进制。
这是整个研究目录的保存成本，后续可无损压缩归档；不将删除负结果作为节省存储的验收方式。

\subsection{第七批：同刻Stream仲裁的增量实验}
当已验证best方案的首批可启动任务总核数或总分配量超出硬件容量时，
尝试逆Stream序、小初始分配优先与高关键路径优先三种编号排列。
保持每条Stream内的任务顺序、原始数据同步及所有额外guard，且保留已有最佳结果；
额外预算最多3次精确评估。这利用题面同刻按Stream编号尝试启动的规则。
第七批压力集153/153、开发143/143、持出143/143、CTest3/3与独立单文件编译全部通过。
相对主体配对版本，压力集$T$改善6例；本轮另行记录的持出增量改善1例，开发集指标不变，均无质量回退。
压力集增量交替测量共同有效方案的solver总时中位1.3694至1.4062秒，变化+2.7\%，
这是小幅CPU成本换取调度质量，不能包装为加速。例$387\to366$、$11371\to11119$、$11024\to10705$。
增量压力CLI最高工作集约449.7MiB，数据文件和前述研究集不同。
原始证据为\texttt{paired-stream-priority-summary.json}与\texttt{development/holdout-comparison-batch7.json}；
该压力配对基线是冻结的\texttt{research/batch6/}，不等于最初的\texttt{baseline-v1}。
增量runner早期将自选基线的源码hash指向v1，已按冻结batch6源码修正并保存provenance说明；原始计时未改。
报告中的主体配对表与此增量表属于不同版本和workload，按各自源/二进制hash归属。

\section{后续潜力与可推进余量}
\subsection{第八批：结构变化、稀缺核心与独立并行oracle}
新增开发种子20261012与持出种子20271012，各72条结构/资源图和24条小图。
其中12条共享零输出relay连接多个显存汇聚分支，打破先前不连通分量的局部性；
36条采用随机长距离边、密度0.02/\allowbreak{}0.12/\allowbreak{}0.35与多sink；24条使用8个Stream但AIC/AIV分别只有8核，
请求1至8核。稀缺核心witness的时间和峰值由第三份Python事件模拟独立计算，再由两个C++模型逐条检查。
全部witness均在运行求解器前构造，数据不根据solver成败筛选。

第七批基线在开发/持出结构集分别61/72与60/72；祖先目标完成遍历回退恢复所有失败，均达72/72。
该迭代DFS先选择低输出目标，完成其全部祖先再处理其他分支，保留原始数据边与直接消费者生命周期。
只在此前策略尚无可行解时追加，复杂度为$O(n\log n+m)$，不会递归耗尽宿主调用栈。
小图则枚举原始DAG所有线性扩展、两流归属与向前的额外跨流事件边；
非循环计划均存在这样的扩展，事件号按记录节点归一，因而在\emph{本地显式直接同步模型}内穷举完整。
独立oracle只链接公共验证器，不调用提交solver；提交文件中另行生成并用自己的模型筛选候选。
规模限定至4节点、两流，不能推导任意大图或官方允许传递同步情况下的全局最优。

\begin{table}[ht]\centering\small
\begin{tabular}{lrrr}\toprule
第八批workload & 第七批/第八批可行数 & 共同图CPU前/后(ms) & 最优T和peak命中\\\midrule
结构开发72条 & 61/72 $\to$ 72/72 & 4152.823 / 4110.117 & 不宣称最优\\
结构持出72条 & 60/72 $\to$ 72/72 & 3431.728 / 3425.505 & 不宣称最优\\
并行oracle开发24条 & 23/24 $\to$ 24/24 & 0.814 / 23.298 & 22/24 $\to$ 24/24\\
并行oracle持出24条 & 21/24 $\to$ 24/24 & 0.782 / 19.068 & 19/24 $\to$ 24/24\\\bottomrule
\end{tabular}
\caption{各集3组交替测量，中位CPU是共同可行图的合计时间；新增可行图不混入该成本分母。}
\end{table}
小图完备搜索成本显著高于原启发式，但整组求解约数十毫秒；不能将该质量收益描述为CPU加速。
结构集共同图计时差在约1\%内，视为基本持平。独立oracle共评估41580个候选（含重复计划），
最终两种子48/48达最优T/peak。新增192条与历史439条构成631条case-record，共享样例重复不计为唯一图。
增量证据位于\path{reports/paired-frontier-batch8-summary.json}、
\path{reports/paired-frontier-holdout-batch8-summary.json}与两个\path{paired-oracle}摘要；
基线冻结在\path{research/batch7/}，本批源哈希\path{df6bec285795bd92}。
主体第六批、第七批仲裁增量与本批工作负载保持分栏，不能互换归属。

\subsection{第九、十批：五节点最优对照与事件剪枝}
独立oracle进一步扩到两流5节点，两个新种子各24图，共穷举1111850个候选（含重复），
生成墙时合计3.171秒；该生成成本与solver测量分开保存。
第八批solver在最优deadline下分别22/24、21/24可行。
第九批先枚举全部无额外Event的拓扑顺序/两流归属，再追加单边Event，额外评估预算6000；
结果两集均24/24最优T/peak，但这一受限五节点搜索不完备，不能推广为一般最优算法。
共同可行案例CPU中位合计分别由1.168至80.057ms、1.362至175.659ms；
这是明确的质量与成本交换，且不同批次的运行负载使绝对计时波动。

第十批缓存精确评估结果并记录是否出现核数/显存阻塞。
同一流内队列若原始计划无资源阻塞，添加Wait/Record只会推迟计算，不能让任意节点早于原始开始时刻；
因此其$T$已经严格差于best时，无需尝试单边Event。保留$T$相等的候选以探索峰值改善，
也保留所有阻塞/非法候选，避免错误排斥通过事件调整竞争的机会。
第九/十批3轮交替配对：开发CPU合计101.623至75.841ms，持出111.338至107.446ms；
精确评估数量分别30049至25803、40860至35215，均约减少14\%。
48/48仍达独立oracle最优，未出现T/同T峰值回退；完整679条case-record回归通过，
CTest3/3与独立单文件编译通过。本批源\path{d4dcbeaacf60f0bc}，冻结对照在\path{research/batch9/}。
新增48个五节点图不应与第八批192条、第六批主表或第七批仲裁表混为同一workload。

\subsection{第十一批：真实消费者窗口的可行性恢复}
新增两种子各30图，规模25至4612节点，多个共享relay连接直接消费者窗口。
每个较大输出由连续2、3或5个消费者使用，只有最后消费者计算结束才能释放；
所有60图有独立构造的串行witness，第十批0/60可行，区别于额外Event不改变释放规则的模型。
新增优先级按四项词典序比较：当前最后消费者的释放量、即将解锁消费者的释放信用、
共享存活输入和最少尚未完成的输入。用lazy版本优先队列维护，邻居更新限制fanin不超过64，
避免密图二次共父遍历成本失控。策略是一般数据边启发式，没有识别固定样例ID或窗口标签。

两种子均30/30恢复；中间将未来信用与共享输入直接相加只恢复9/60，
随后加入小输入窗口仍仅20/60，重复扇出信用会诱导过多半完成窗口，
拆为独立词典序后60/60。更关键的是，先把新策略插在旧仲裁前导致1条压力图$2361\to2378$回退，
再仅在早期nofound触发仍产生旧局部性回退图的等T峰值回退。
最终将其严格放在\emph{全部已有搜索及Stream仲裁之后，且仍未找到可行候选}才运行，
保留此前所有可行输入的完整轨迹。被拒绝草案源hash\path{08fc39e010428317}保存于
\path{research/batch11-rejected/}，负结果与输入均保留。

第十一批最终源\path{5a8518a47ea18852}完整739条case-record回归，所有质量指标无回退，
独立模型一致；新增60图没有共同可行基线，\emph{不能计算共同T改善或将共同CPU=0称为加速}。
仅报告完整CLI流程成本：开发失败基线/成功新版中位1.028/\allowbreak{}1.042秒，
持出0.963/\allowbreak{}1.006秒。对应的完整搜索成本既含旧失败候选也含新增恢复，
此比较不具备同质量CPU加速含义。全739图功能回归墙时24.490秒。
固定均匀输出的窗口族仍不能代表任意权重、任意长距离共享消费；后续批次再扩大异构权重与前沿搜索。

\subsection{第十二批：大图资源调度质量与明确CPU代价}
新建两种子20261016/20271016各18个稀缺核争抢图，节点数64/256/1024，
边密度0.03/\allowbreak{}0.1，Event单价0/3/$10^5$。witness仍由独立Python资源模型生成，
对照第十一批已有36/36可行，因此这里评价共同可行图的质量，而非恢复率。
在所有旧搜索与仲裁之后，追加最多4个关键路径/释放优先的保守资源列表候选，
Stream数4/8并分开有无完成前沿guard，完整模拟只接受严格更优$(T,peak)$。
这样不会替代原本已检查的最佳方案，也不按case标签选择策略。
开发9/18、持出8/18的$T$降低；最大分别31.17\%和28.02\%，其余共同图没有质量回退。
三轮交替测量共同solver CPU合计中位为395.186至443.061ms、
407.235至453.873ms，成本约增加12.1\%和11.5\%。
这是通用大图调度质量收益，需要额外CPU，不是计算加速；本批源\path{19a9686e14c07c8c}。

\subsection{第十三批：非均匀内存前沿与有界beam}
将第十一批窗口族的输出改为60至140、工作区150至300、消费者时长1至50、sink输出0至3，
两新种子20261015/20271015各30图，规模25至4612节点。所有图保留独立串行witness；
第十二批均0/30可行，说明均匀内存优先级无法代表异构权重需求。
新增分类前沿策略将当前真正释放、已有共享输入、打开新消费窗口分成三个优先层级，
只有最后消费者贡献实际释放；尚未就绪消费者的释放信用仅影响排序，不改变生命周期。
仍在全部旧搜索后尚无解时触发，最大节点数10000。

512节点以内还搜索宽度8、每状态最多6分支的beam。状态记录已完成集合、剩余入度/消费者计数、
ready集合、存活内存与历史峰值；仅展开串行分配可容纳的ready任务。
按已释放总量、活跃消费者剩余输入数、当前存活量和峰值排序，
相同完成集合使用哈希加完整向量检查合并，保留较低峰值。
不同顺序的后续串行内存状态由集合确定，但这种合并不保留所有潜在并行分流顺序，不能推导并行最优。
beam与分类贪心形成候选，再尝试串行和少量分流，最终仍由完整精确模型筛选。
两集均30/30恢复；36个小图启动beam，共扩展194864个子状态，最大保留8个状态。
其余24个超过512节点的图依赖分类贪心候选，不使用beam；
这些统计仅表明beam被尝试，不能将组合的全部恢复收益归因于beam。

新增60图无共同可行基线，不能报告共同$T$改善或CPU加速。
完整CLI失败旧版/成功新版中位开发1.316/\allowbreak{}2.350秒，
持出1.075/\allowbreak{}1.888秒；后台负载导致单轮持出墙时最高约9.95秒，故保留全部原始重复数据。
新版835条case-record完整回归墙时29.989秒，独立模型一致，无$T$/同$T$峰值回退；
CTest3/3用时0.56秒，独立单文件编译无警告。当前源码\path{5faa403684db1243}，
对应两个\path{paired-weighted-beam}摘要及\path{reports/regression-batch13.json}。

\subsection{第十四批：不规则长距离共享、部分恢复与成本控制}
新增两种子20261017/20271017各24图，阶段内生产者1至3、伴随输出1至2、消费者1至3均随机，
消费关系采用非连续的局部跨度8、密集跨度32或全图跨度；可选消费者小输出的后续复用。
各阶段只用于构造独立serial witness，不添加强制阶段顺序边，ID与记录乱序，图不按solver成功筛选。
与前面固定窗口族相比，规模42至2151、核请求1至8、不同fanout与复用输出造成不同存活割。
第十三批仅开发3/24、持出6/24，证明已有60条窗口全通过不足以代表一般共享消费。

新增目标束前瞻：针对一个尚未完成的目标，用迭代DFS取其缺失祖先，最多64节点，
临时按真实直接消费者计数模拟整个串行束的分配/释放，探测后回滚计数；
只有各步骤都能容纳才可选择。分别试净释放、最少新输出、既有输入信用/新输出比率三种排序，
总额外候选预算9，图规模至4096。目标束完成后重新评价，不把输入信用冒充已释放显存。
原逐ready任务优先级会打开多个半完成消费前沿，而整束能够提前发现消费者工作区不足。
该策略先将9/48恢复到20/48。额外尝试完成某存活生产者的\emph{全部剩余消费者}未恢复更多图，
且增加扫描成本，因此没有保留；探索性草案只有原始结果，不具有冻结源码的配对计时归属。

未解的512节点以内图继续尝试宽度32、分支12的两种beam排序：
原释放量排序，以及按尚未完成消费者比例折扣的存活输出信用，仍检查完整实际显存。
再恢复3图，最终开发11/24、持出12/24，48条中仍25条\emph{已知有可行witness但搜索失败}。
错误和完整失败case保存于各paired的\path{.failures.json}，不把失败改标为expected-infeasible。
原835条完整回归通过，墙时32.050秒；CTest3/3用时0.93秒，单文件编译无警告。
总883条记录中858条符合测试预期、25条失败；共享官方例重复，不能称883个独立问题或883条全通过。

原版共同可行3/6条的CPU中位合计55.226/\allowbreak{}54.949ms、
75.950/\allowbreak{}75.967ms，基本持平且质量无回退；整CLI因恢复搜索增加，
开发0.875至2.595秒、持出0.837至2.651秒。
对超过512节点的束前瞻，再以四种廉价估计各取前32目标、并成最多128项，避免每轮探测全图目标。
与冻结的同算法全扫描版本比较，成功数和共同T/peak相同，
共同CPU开发279.124至220.760ms、持出1654.598至924.018ms，
分别降低20.9\%和44.2\%；这是扫描实现成本改善，与相对第十三批的恢复成本不是同一分母。
冻结全扫描源/二进制位于\path{research/batch14-full-scan/}，当前源\path{29eeb08516473d15}。
有界bundle/beam仍无法覆盖随机非连续共享消费，未来价值在更强回溯、状态下界和可行性证据利用。

\subsection{第十五批：六节点最优标尺与同步搜索}
独立oracle扩到六节点两流，两种子20261018/20271018各12图，
枚举所有原始线性扩展、流归属和向前跨流非数据Event边子集，共11996268个候选，
生成墙时38.300秒与solver时延分开记录。最多6节点时两流跨边至多9，子集编码不会溢出。
独立oracle不调用提交solver，完整性仍只对应本地显式直接同步模型。
第十四批最优deadline下开发11/12、持出5/12可行，最优T/peak分别仅9/12、4/12。

提交侧六节点首先枚举无额外Event的顺序/流归属，再分配12000次单边和8000次双边同步评估，
各阶段独立预算，达到预算即停止无作用循环；沿用无资源阻塞且T严格更差的事件剪枝。
双边候选不是任意Event子集，提交侧仍不完备。
最终两集12/12可行且24/24 T/peak命中独立最优，恢复8条并改善若干已有可行图的峰值。
代价明确：共同11条CPU由1.255至774.325ms，
共同5条由0.341至213.625ms；不能描述成CPU加速。
完整859条严格回归通过（56.682秒），第十四批困难图仍23/48，无既有成功图质量回退。
907条总记录882通过、25条有witness但未解，CTest3/3与单文件编译通过。
本批源码\path{9e4a633813a59b08}与\path{research/batch15/}冻结一致。

\subsection{第十六批：sink祖先闭包投影回溯}
剩余困难图中，有46节点9个sink、48节点14个sink的例子，完整节点子集DP已不可用，
但最终输出数量较少，因此以完成sink作为投影状态。
预计算每个sink的祖先闭包及固定迭代DFS顺序；选中sink集合$S$时，
完成节点集合恰为这些闭包的并集$C(S)$。原始输出存活量与直接消费者剩余数仅由$C(S)$决定，
与之前完成sink的顺序无关；已经完整探索失败的$S$可缓存剪掉其他到达路径。
转移执行目标sink尚缺的祖先，逐步检查实际workspace+output和最后直接消费者释放，
按转移后存活量与局部峰值排序后回溯。
上限512节点、18 sink、800万次元素操作；超过预算返回无候选，不能据此证明无解。
固定闭包内顺序未枚举所有节点拓扑序，所以这不是一般串行最优算法。

原48困难图23/48至25/48，额外恢复上述两图；以新种子20261019/20271019、6/8阶段
再构造各12条不规则消费图，原10/12、11/12均至12/12，共新增恢复5条。
额外新图按输入规则直接生成，未按sink数量或solver成败筛选。
原859记录完整回归用时38.231秒，另24新小图在交替配对中全部通过，
所有已成功图质量无回退，独立模型一致。CTest3/3用时0.71秒，单文件编译无警告。
当前931条总记录908通过、23条有独立witness但未找到方案；仍未覆盖更多sink的大图。
两套小图共同CPU合计开发44.144/\allowbreak{}45.690ms、
持出51.465/\allowbreak{}51.197ms，旧成功路径未调用新回溯，微小变化应视为后台计时波动。
本批源\path{6921067fb0f8c567}，冻结于\path{research/batch16/}；
859条单次回归和另24条配对验证是不同测量，不合成虚构的883条单次回归墙时。

\subsection{第十七批：按共享输出分组扩大投影范围}
多sink使18位投影无法直接使用，因此进一步将最多72个sink按共享正输出父节点分组。
父节点权重取输出量除以扇出数，限制扇出不超过16，降低反复计算高扇出输入信用的影响；
按共享权重降序合并，每组最多4成员，直到18组。不能合并到18组便返回无候选。
每组的祖先闭包仍由原始数据边构造，临时信用不释放显存；
组内固定DFS顺序和组间有界回溯只搜索部分合法计划，仍没有一般可行性或最优保证。
旧18 sink以内的路径保持一致，新分组只处理此前超过18 sink的未解图。

原48困难图25/48至29/48，开发12/24至15/24、持出13/24至14/24；
已有成功图T/同T峰值无回退，883条严格回归全部通过（48.856秒），
CTest3/3用时1.60秒、单文件独立编译无警告。931总记录912通过、19有witness仍未解。
首个草案未排除没有后继的正输出节点，亲和权重除零使CLI运行失败；
修复该条件并重建后才执行正式配对，不使用那次没有有效输出的运行作成功或成本证据。
共同12/13条CPU合计中位分别270.375/\allowbreak{}308.069ms、
1241.292/\allowbreak{}1069.006ms；旧成功路径未调用新分组，差异受并发后台负载影响，不解释为稳定加速。
完整开发CLI中位3.460/\allowbreak{}3.137秒、持出3.224/\allowbreak{}3.378秒，
亦不将失败恢复的流程成本包装成共同质量加速。
本批源\path{87cade82e33518fe}冻结于\path{research/batch17/}，各paired完整失败输入仍保存。

\subsection{第十八批：核需求优先级与跨分布收益差异}
将研究重点回到通用图完工时间。保留第十二批4个资源列表候选，再独立追加4个候选：
按对应核池容量/请求核数归一关键路径秩，以及按去除自身时长后的剩余关键尾排序。
两种排序都用真正的最后消费者释放作为次级项，Stream数4/8，无额外前沿guard。
所有旧搜索及候选已完整执行，精确评估仅接受严格更优T/peak，因此新策略不抢占旧预算。

已有36条资源DAG共同可行，开发6/18、持出7/18进一步降低T；最大开发12.33\%，
部分小改善仅数千微秒，不能将改善数量视作同等幅度收益。
共同CPU开发494.762至533.817ms、持出442.042至528.278ms，
增加约7.9\%和19.5\%，仍为质量换取更多搜索工作。
额外两种子20261020/20271020各18条核装箱图：64/256/1024节点、8/32核池，
无数据依赖，时长分别独立于核数、按核数反比例缩放、按核数正比例缩放；内存上界容纳任意8个并发计算，
每图具有独立serial witness。新分布36/36共同可行，但仅持出1条T降低2.92\%，开发3条同T峰值改善。
这表明新优先级对有依赖的核争抢图更有效，不能宣称泛化为普遍装箱优势。

核装箱三轮中位CPU开发625.135/\allowbreak{}190.060ms、持出1041.160/\allowbreak{}455.162ms，
前后并发后台负载显著不同；源码追加了4候选，并没有删除旧工作，
这些表面计时下降不归因于算法降低计算成本，也不称为稳定加速。
完整919条严格回归用时34.987秒；独立模型一致，无旧T/同T峰值回退，
原48困难图仍29成功19失败。967总记录948通过19失败，CTest3/3用时0.71秒、单文件编译无警告。
本批源\path{6e66b643ef5a58d8}冻结\path{research/batch18/}；不同工作负载的T收益与CPU成本分别保存。

\subsection{第十九批：Wait与Compute资源分离的构造器}
原资源列表在Wait阶段预先保留核心与内存，会阻止其他可计算任务利用等待窗口。
新构造器区分pending Wait与active Compute，只在其计算阶段真正启动时分配资源；
原始直接消费者结束时释放输出，按Stream编号尝试启动，资源不足时继续其他流。
队列构造中仍给所有生产者结束保守计入Record，尚不知道最终哪些输出会有跨流消费者，
故其虚拟时间不宣称与最终提交计划完全等价；发出的每个计划须经过独立的精确评估器确认T/peak与合法性。
只在Event耗时非零时追加至多4个质量候选，避免零Event图的额外无收益构造。
未解图另试4个恢复候选，但本批没有恢复原困难集，明确保留该负结果。

旧36条资源DAG只有2条T改善，幅度0.26\%和0.33\%。为验证适用范围，
独立生成两种子20261021/20271021各18图：2/4/8个独立宽汇聚分支，
每分支8/32/128个16或32核生产者，汇聚消费者要求满128核，Wait单价100/$10^4$，Record取其一半。
18至1032节点，内存保守容纳任意8个并发计算，所有图有独立serial witness。
新分布开发8/18、持出7/18的T降低，最大6.38\%；收益全部在Wait100组，
Wait$10^4$组没有改善，不能推广为越大Event越有效。
构造器只有部分跨流同步时延能与计算重叠，极高Event下已有串行或低同步计划可能更好。

共同资源DAG CPU开发502.371至549.308ms、持出481.327至527.298ms，
约增加9.3\%/9.6\%。新汇聚共同CPU开发177.580至222.174ms、
持出157.567至210.742ms，约增加25.1\%/33.7\%；这些是明确搜索成本。
955条严格回归通过（60.582秒），独立模型一致、无T/同T峰值回退，
原困难48图仍29成功19失败；1003条总记录984通过19失败。
CTest3/3在并发负载下用时10.95秒，单文件编译无警告。
另将核归一优先级改为整数商/余数分解并饱和，避免较大核池数乘法溢出；本地合法数据排序数值保持相同。
补充的4项边界验证覆盖结果仍可表示但直接乘法溢出、超界饱和、不能整除时的取整，独立常量期望全部通过。
本批源\path{efa449b6bcef7656}冻结\path{research/batch19/}。

\subsection{第二十批：最终Stream编号试探的收益与默认取舍}
在现有最优候选确实发生过核心或显存阻塞、节点不超过2048时，额外试反转及三次旋转Stream编号，
保留队列、原始数据边和已有额外Event，使用独立预算4次精确模拟。
该结构依据是同刻启动按Stream编号决定，阻塞可出现在图中间，原有起点排序不能覆盖全部。
108条资源DAG、宽汇聚和核装箱记录中，仅两条T降低（约0.001\%与2.04\%），另有一条同T峰值降低。
宽汇聚持出18条配对CPU为208.634至222.116ms，增加6.5\%；全CLI为0.398至0.427秒。
资源持出CPU586.701至519.237ms波动下降，但源代码增加搜索，没有减少旧计算，且三次重复明显受共享负载影响，
不把这项波动当算法降本。
955条严格回归44.442秒，无非法返回、模型差异或相对于既有基线的质量回退；困难集仍29/48。
完整源\path{0d2a1bd218f5744f}及原始结果冻结\path{research/batch20/}。
因为收益范围窄且存在普遍额外评估成本，默认求解器不保留本批四次重编号，后续以第十九批为质量对照。
这明确放弃上述少量探索性质量收益，不能声称第二十一批相对于第二十批所有案例质量无损。

\subsection{第二十一批：已证最优候选的搜索终止证书}
本批不增加候选，而用必要下界判断已经无可改进的计划。记算子计算耗时$d_i$、核数$c_i$，
Stream数量$S$、两类核池容量$C_t$，时间下界为
\[
L_T=\max\left\{\text{最长原始数据路径},\left\lceil\frac{\sum_i d_i}{S}\right\rceil,
\max_t\left\lceil\frac{\sum_{i\in t}d_i c_i}{C_t}\right\rceil\right\}.
\]
前两项分别来自先后依赖和Stream内计算不重叠，核工作量来自计算区间内资源积分上限。
Event成本忽略，故只会使下界更弱；空核池只有对应无任务时允许，累加饱和取小不会高估真实必要量。
算子启动必须同时持有自身workspace、output及所有直接父输出。因此显存下界包含
\[
L_M\ge\max_i\bigl(W_i+O_i+\sum_{p\in pred(i)}O_p\bigr).
\]
所有sink输出持续到全图结束。考虑最后启动的sink，可进一步加入
\[
L_M\ge \sum_{s\in sink}O_s+\min_{s\in sink}\bigl(W_s+\sum_{p\in pred(s)}O_p\bigr).
\]
直接父节点不是sink，二式没有重复计数。

当候选$T=L_T$且$n\le128$，根据原始DAG仅计算耗时求最早启动$E_i$与最晚结束
$F_i=L_T-(rank_i-d_i)$。若先$u$后$v$的必要条件
$\max(E_v,E_u+d_u)+d_v\le F_v$及反向条件都不满足，则在该最短时间内两算子必然有正长度计算重叠。
其两个workspace/output及去重后的直接父输出之和构成\emph{条件}显存下界。
窗口忽略Event和资源，因此仍包含全部真实计划；两个非重叠顺序皆被排除才能作该推断。
只在已验证候选同时达到时间和显存下界时提前返回，不依赖候选优先级的经验推断。
条件显存界不用于比较更慢计划的最优peak，也不用于提前宣布无解。
证书适用于题面固定资源模型，不等于真实NPU运行性能保证；一般困难图仍是有界启发式。

独立解析控制包含两种子的1/2/32/512/10000节点总链，以及两个同核池或异构满核任务，
共12条，随机稀疏ID与乱序记录，链耗时达64位范围。总链只有一种数据拓扑顺序，
任何合法计划均不能并发计算，最优T与peak可独立计算；异构两任务在最优T内必重叠，peak=4。
已冻结的3至6节点完整并行oracle另提供120条对照，不能把witness代替最优标尺。
132个独立最优标尺检查通过，其中21个有证书；955条广泛回归有86个证书，12条解析控制全部有证书。
严格回归60.685秒，困难图仍29/48；加解析控制后1015条记录996通过19有解仍未找到计划。
相对于默认第十九批共同案例的T与同T峰值完全保持。
三次交替配对，解析控制CPU25.066至8.092ms，降低67.7\%；
主开发143图2892.121至2868.459ms、持出143图2997.631至2964.331ms，只有约0.8\%/1.1\%下降，
共享负载下不能把这些小幅变化当作普遍稳定加速。
无证书命中的6节点oracle持出12图CPU640.289至670.465ms，增加4.7\%，保留其负成本对照。
解析控制全CLI0.648至0.404秒，但JSON/验证/共享干扰占比大，求解部分单独列出。
CTest3/3用时0.70秒，单文件GCC独立编译无警告；本批源\path{f95ea2f74d1d37b2}冻结\path{research/batch21/}。

\subsection{第二十二批：密集输入的等价分流代数优化}
原分流递推对每算子、每目标Stream再次遍历所有直接父节点，成本$O(S(n+m))$。
现在先按父所属Stream汇总最大完成时刻与数量，记录最大、次大组的完成时刻，
则对目标$s$，同流父节点最晚时刻可直接读取，跨流父最晚时刻等于排除$s$组后的组最大值再加Record。
跨流Wait数量等于总父数减同流父数；并列组最大值通过次大值正确保留。
保持原有affinity惩罚与最早Stream打破并列规则，在fanin至少4、Stream至少4时使用汇总，复杂度$O(m+nS)$。
另将初始排序按需生成：首个串行候选已有双界证书时，跳过剩余排序与串行候选。

以冻结的逐父递推作独立对照，288组随机参数覆盖并列完成、Event0/1/$10^6$、64位时间、
乱序ID/记录与1/4/8/16个Stream，分流数组全部相同。
1015条既有/解析/困难记录的最终计划JSON及合法性、T、peak逐条相同，
955严格回归39.255秒、132独立最优标尺审计通过，原困难仍29/48。
新增两种子20261023/20271023各6个密集图，节点256/1024/4096、层宽16/32/64、fanin8/24/48，
Stream16/64/256，分别含相邻层与任意此前层输入；没有总链，不按solver结果筛选，均有独立serial witness。
12输入文件合计约12.8MB，最大全CLI工作集127.2MiB，不是设备张量或提交函数独占内存。

共同案例T/peak与返回计划保持；solver CPU开发2382.426至2076.588ms、持出2463.341至1843.333ms，
分别下降12.8\%/25.2\%；全CLI开发5.893至4.680秒、持出4.049至3.137秒。
泛化主持出143图CPU3072.623至3727.082ms增加21.3\%，同时全CLI5.679至5.603秒，
并发宿主负载有明显干扰，仍保留负成本记录，不能把密集输入的收益推广所有DAG。
解析控制CPU7.791至7.229ms，额外按需排序只有小幅下降；其全CLI反而1.487至2.978秒。
后续应提高汇总触发阈值并延迟scratch分配，防止小fanin的汇总开销。
新12图全部通过后1027条总记录1008通过19有witness仍未解；CTest4/4用时0.80秒，单文件独立编译无警告。
本批源\path{ef3797e37fe6c776}冻结\path{research/batch22/}，配对与逐计划一致性证据单独保存。

\subsection{第二十三批：局部性负结果、触发阈值与计时口径审计}
尝试在仅用于优先级的无向图中过滤零输出边，或进一步过滤高扇出父节点的边；
全量原始数据边仍通过Kahn排序、Event和输出释放保留，未修改真正依赖。
结构依据是大量零输出共享根会把存储阶段连接为星状，掩盖较局部的正输出生命周期。
但在原48困难图没有额外恢复；新增两种子20261024/20271024，阶段数12/32/96、
非连续局部/密集/全图共享、可选消费者小输出后续复用，以及共享根输出0或5，共72图，
有独立serial witness且没有按solver成败筛选。
开发旧/新均22/36、持出均21/36，新增29个有解未找到计划的覆盖缺口明确保存。
追加搜索造成失败路径开销，草案冻结\path{research/batch23-rejected-lifetime/}后从默认撤回；
这些单次诊断只支持无恢复结论，不作为稳定CPU增幅的正式配对数值。

默认改动只提高汇总触发门槛至fanin至少8、Stream至少4且二者乘积至少64，
scratch到首次真正汇总时才分配。其依据是小fanin/少Stream时汇总、触摸计数及分支开销可能抵消遍历收益。
属于同一分流递推的等价变换，不引入新调度候选。
979条严格回归65.803秒；1027条既有计划逐条相同，新增72条输出计划/成败亦与冻结22相同，
132个独立最优标尺审计通过。1099总记录1051通过48有解未找到结果，其中旧困难19、新结构29。
三次交替配对主持出143图求解墙钟2682.432至2728.553ms（增加1.7\%），
线程CPU2656.250至2718.750ms（增加2.4\%）；密集持出6图墙钟1785.155至1786.866ms近乎不变，
线程CPU1765.625至1734.375ms下降1.8\%。没有观测到可观阈值加速，保留为小fanin的成本触发保护，
不将量化/共享资源范围内的小幅变化作为收益结论。CTest4/4 5.09秒、独立编译无警告。
本批源\path{4c5b2c6a32b72331}冻结\path{research/batch23/}。

审计发现历史\texttt{solve\_us}是CPU求解器\emph{墙钟}，包含OS抢占，而非线程实际CPU。
本地harness新增\texttt{solve\_thread\_cpu\_us}，通过Windows \texttt{GetThreadTimes}读取kernel+user之差\cite{threadtime}；
不支持的平台记$-1$，从未将不可用当作0。该Windows依赖只在本地benchmark，提交源仍仅problem.h与标准库。
官方FILETIME单位100ns不是本机实际更新时间粒度；实测计数按15625us步进，单图0不能解释成没有CPU代价。
大量小图聚合也存在量化偏差，线程CPU作为辅助口径，墙钟/全CLI继续单独报告。
为了公平比较，冻结第二十二批提交源使用完全相同的新harness重编，
源/二进制/harness/simulator/CMakeCache哈希记录于\path{research/batch22-instrumented/provenance.json}。
旧历史数字不套用新时钟，也不因口径审计改写。

\subsection{第二十四批：稀疏谱平滑的有限恢复与拒绝默认}
Fiedler经典工作\cite{fiedler}使用图Laplacian的第二特征值描述代数连通性，
为本地通过低频图表示整理消费局部性的尝试提供背景，不能直接推出题面内存最优性。
本批建立正输出父边的无向加权图，边权为父输出/直接消费者数量，忽略扇出超过32的优先级hub；
原始数据边、Event和输出释放不变。每连通分量去除degree加权均值并归一化，
进行256轮惰性邻居平均，使用BFS远端距离或确定性ID散列初始化。
固定迭代数只得到受谱方法启发的坐标，不声明已求准特征向量。
坐标正反方向分别生成Kahn优先顺序和目标祖先DFS完成顺序，未解且$n\le4096$时追加至多40候选。

120条原困难/新结构单次定位中，原48仍29成功；新72由43至44成功，
仅持出密集12阶段、保留小输出且共享根输出5的一图恢复。
随后将sink或workspace大于等于output的目标提前完成，也没有额外恢复。
如此窄收益不足以支持普遍256轮平滑与40次追加搜索，本批冻结\path{research/batch24-rejected-spectral/}后撤回默认。
独立模拟验证恢复方案合法，但没有进行无收益草案的全面回归或正式配对成本测量，
不能把单次运行时差当稳定增幅或宣称一般图可行率提高。
本批默认保留第二十三批；后续第二十六批再以较直接的消费组束替代无收益的全局排序追加。

\subsection{第二十五批：祖先次序与近优束重启的零恢复}
先尝试按父输出从小到大或平均释放信用排序目标祖先，各配合三个原束评分，共六种次序/评分组合。
仍使用原始直接用户计数精确检查整个串行束的临时分配和释放，越界探测回滚。
120条定位中没有额外恢复，原困难29/48、新结构43/72保持。
再限制$n\le768$，以固定种子进行12次重启：保留评分最好的六个合法目标束，
约一半选择最好目标，其余在近优集合中选择；每次仍逐节点检查实际内存，完整计划必须精确模拟。
此多起点试探同样没有恢复120条中的任何新案例，不能将随机多样性当作实际收益。
两版冻结\path{research/batch25-rejected-ancestor/}和\path{research/batch25-rejected-restarts/}后全部撤回默认。
仅进行了聚焦质量诊断，没有以无收益草案反复做全量回归或正式配对成本测量。
固定单消费者目标不能主动形成较完整的共享输入释放单元，下一步将探索小组消费束，
保持原图释放规则且避免混入错误的额外依赖。

\subsection{第二十六批：选择性共享输入消费组束}
单个目标消费者的祖先束可能打开共享输入，却把其余消费者留在后续前沿，延长持有。
新方案只用于原候选全部失败且$n\le4096$的图；不会为已成功图追加质量搜索。
在正输出且直接消费者数为2至4的父节点上，选择sink或workspace大于等于output的子节点，
并要求父执行时长乘所选消费者数不大于消费者时长之和，作为相对短生命周期输入的触发保护。
以输出量/扇出/执行时长形成消费对权重，排序后用并查集合并，组上限4；未分组节点保留单体。
按组完成祖先束，同组优先处理大workspace消费者，临时分配及释放始终使用原图直接后继计数，
不添加数据依赖，不提前释放输出，也不把优先级组当作题面执行原子。
束祖先上限由64至128，三种原评分各尝试串行及4/8流，去重后至多9候选，首个通过即停止。

先试不含相对执行时长保护的原组束，仅恢复原困难的一图；保护版本合计恢复5条。
正确的消融保留同样128祖先上限、评分、候选预算和流数，但禁用消费分组，
对选出的48条此前失败记录恢复0；这是定位子集证据，不能推广为分组对任意DAG有效的证明。
一次误置消融开关实际触发旧束，结果已单独重命名并排除；一次冻结源复制保留旧mtime导致CMake复用对象，
在完整报告形成前停止，触摸源后确认编译/链接再运行以下正式流水线，未混用错误二进制。

正式979条严格回归49.719秒，无T退化、同T峰值退化或模型分歧；132独立最优标尺审计通过，21个双界证书。
原困难48由29至30成功，恢复2095节点全局共享图；新生命周期72由43至47成功，
恢复局部32阶段两种根输出控制、局部96阶段根输出0，以及密集12阶段根输出5。
两个32阶段控制共享基础图，只计作不同参数case-record，不能当两种独立拓扑。
1099总记录1056通过、43个有witness但未解；CTest4/4用时3.37秒，单文件独立GCC编译无警告。
共同可行T/peak全部保持，恢复可行不能当共同案例加速。源\path{a72873d1ab827d51}冻结\path{research/batch26/}。

以下同一harness/冻结23与26二进制，各做三次交替配对，中位数；求解成本包含成功与失败的所有记录。
\par\smallskip\noindent{\small\setlength{\tabcolsep}{3pt}\begin{tabular}{>{\raggedright\arraybackslash}p{3.2cm}>{\raggedleft\arraybackslash}p{1.5cm}>{\raggedleft\arraybackslash}p{3.1cm}>{\raggedleft\arraybackslash}p{3.1cm}>{\raggedleft\arraybackslash}p{1.8cm}}
集合 & 可行数 & 调用墙时ms & 线程CPU ms & CLI秒\\\hline
原困难开发24 & 15/16 & 2671.398/\allowbreak{}3500.128 & 2546.875/\allowbreak{}3421.875 & 3.733/\allowbreak{}4.276\\
原困难第二种子24 & 14/14 & 2387.833/\allowbreak{}3019.832 & 2296.875/\allowbreak{}2734.375 & 3.118/\allowbreak{}3.304\\
生命周期开发36 & 22/25 & 1404.216/\allowbreak{}1509.604 & 1375.000/\allowbreak{}1531.250 & 1.714/\allowbreak{}1.770\\
生命周期第二种子36 & 21/22 & 1417.248/\allowbreak{}1955.693 & 1421.875/\allowbreak{}1593.750 & 1.828/\allowbreak{}2.362\\
\end{tabular}}\par\smallskip

相应共同案例求解墙时分别428.874/\allowbreak{}479.008、1206.593/\allowbreak{}1064.337、482.498/\allowbreak{}415.545、421.421/\allowbreak{}495.288ms；
线程CPU分别406.250/\allowbreak{}437.500、1109.375/\allowbreak{}1062.500、406.250/\allowbreak{}406.250、421.875/\allowbreak{}421.875ms。
已成功图不触发新束，变化包含共享宿主抢占及CPU计数量化；不能解释成既有算法被加速或确定变慢。
含失败搜索的线程CPU增加34.4\%/19.0\%/11.4\%/12.1\%，与仅追加9候选仍有实质成本一致，
因此保留未解路径触发及4096规模上限，不对全部图普遍启用。
开发与第二种子在历次迭代均被查看，称为两个固定种子而非严格盲测。
冻结26源码和二进制后，以第三种子20281026生成一次未见36图，不按solver结果筛选，节点60至1292。
每个阶段数12/64/128，局部/密集连接和保留输出0/1组合，分别采用原不均匀时长、
生产者与消费者时长均1至100以消除数量级分离、消费者增至5至8以改变高扇出结构。
输入生成即记录哈希，再以冻结23/26各测一次；未据该结果再调整26，单次成本不是正式配对结论。
原不均匀组3/12至5/12，恢复64阶段局部和密集两图；接近时长4/12、高扇出1/12保持，
总8/36至10/36，共同8图T/peak保持且无模型分歧。
全调用墙时2746.065至3226.589ms，线程CPU2546.875至3156.250ms，仅诊断观察。
独立serial witness和提交模型对所有输入核对通过，26个有解仍未解记录可重放。
合并后1135条记录1066通过、69个有witness未解；高扇出和相近时长的零增益明确保留，
不将5条历史恢复或2条新恢复推断为一般图成功率。

\subsection{第二十七批：主要消费组的遍历先后偏好}
消费组束仍可能先完成未来孤立消费者的祖先，从而较早打开其共享输入。
在第二十六批组上，对每个正输出父节点统计直接消费者所属组，
仅唯一最大组包含至少2个用户时，将该组用户优先于其他孤立/较小组用户，构成\emph{排序专用}先后边。
所有排序边去重后整体Kahn判环，若存在环则放弃本候选；没有将偏好当真实数据依赖或新增Event。
原图分配/释放计数不变，完整计划经精确模型和独立模拟器检查。
只有旧方法及第二十六批全部失败且$n\le4096$才追加三评分的串行/4/8流，至多9唯一候选。
主要组最多4个成员，按每个父节点先收集owner用户后批量追加，避免对高扇出父节点作平方扫描；
排序边空间至多4倍原边规模，Kahn为线性扫描，束评分仍沿用原预算和128祖先上限。

新增6项明确构造检查：三评分各覆盖无环主要组优先、原始拓扑保持及原DAG无环但排序偏好成环的拒绝。
与已有288分流递推检查一同通过；979严格回归49.110秒、132独立最优标尺/21证书审计通过。
原48困难由30至32，新72由47至54，分别恢复两图和七条参数记录，
其中根输出0/5对仍不是独立拓扑。第二十六批36结构图10/36保持，此集已被观察，未称为本批盲测。
累计1135条记录1075通过60个有witness未解；共同T/peak保持，CTest4/4 0.82秒、独立单文件编译无警告。
本批源\path{77ef67341e91a00a}冻结\path{research/batch27/}。

三次交替冻结26/27配对，以下全调用成本包括失败搜索，不能用共同成功分母掩盖追加失败路径：
\par\smallskip\noindent{\small\setlength{\tabcolsep}{3pt}\begin{tabular}{>{\raggedright\arraybackslash}p{3.2cm}>{\raggedleft\arraybackslash}p{1.5cm}>{\raggedleft\arraybackslash}p{3.1cm}>{\raggedleft\arraybackslash}p{3.1cm}>{\raggedleft\arraybackslash}p{1.8cm}}
集合 & 可行数 & 调用墙时ms & 线程CPU ms & CLI秒\\\hline
原困难开发24 & 16/17 & 3143.401/\allowbreak{}4608.843 & 3078.125/\allowbreak{}3828.125 & 4.559/\allowbreak{}5.439\\
原困难第二种子24 & 14/15 & 2572.969/\allowbreak{}3531.555 & 2562.500/\allowbreak{}3312.500 & 3.102/\allowbreak{}5.088\\
生命周期开发36 & 25/27 & 1389.416/\allowbreak{}1562.432 & 1359.375/\allowbreak{}1562.500 & 2.085/\allowbreak{}1.937\\
生命周期第二种子36 & 22/27 & 2806.816/\allowbreak{}2281.090 & 1859.375/\allowbreak{}1781.250 & 3.032/\allowbreak{}2.478\\
\end{tabular}}\par\smallskip

共同案例墙时/线程CPU依次为1403.254/\allowbreak{}1759.510/\allowbreak{}1390.625/\allowbreak{}1625.000、1019.003/\allowbreak{}1060.648/\allowbreak{}1015.625/\allowbreak{}953.125、659.110/\allowbreak{}652.544/\allowbreak{}671.875/\allowbreak{}640.625、700.791/\allowbreak{}587.490/\allowbreak{}453.125/\allowbreak{}453.125（ms，前/后）。
已成功图不触发该方法，不能将其共同计时差解释成算法质量收益；共享宿主负载、15625us计数量化仍存在。
原困难开发某次全CLI16.509秒，相邻冻结26为8.029秒，表中使用各三次中位数且保留原始值。
全记录线程CPU前三组增加24.4\%/29.3\%/14.9\%，生命周期第二种子反而下降4.2\%；
追加方法以前的搜索没有删除，该下降不能作为确定的计算量减少，保留为共享负载/量化下观测。
高扇出/相近执行时长组无增益，限制默认触发，后续应尝试按活跃前沿而非不断扩大所有图的搜索预算。

另评估小图精确串行peak下界扩展：任意合法并行计划按Compute启动顺序串行化，
此前已启动任务在串行中均已完成；串行仍保留的输出在并行中不可能已释放，因为至少一个直接用户尚未启动，
故该串行化peak不大于原并行peak。最小串行peak因而可作本题模型的全局必要界，
不能转用输入在消费者启动时删除的文献模型。
然而当前主两种子共44个小图未命中简易peak证书，其返回peak已等于独立serial oracle，
求解总墙时仅9.152ms；重新作$n2^n$DP仅增加证书数量，暂无可观质量/成本收益依据，未加入默认。

\subsection{第二十八批：宽消费组及结构触发保护}
针对高扇出组，将短输入的有效直接消费范围、合并组上限由4扩至8，
保留相对执行时长条件、主要组遍历偏好和真实直接消费释放。
仍只在此前方法全部失败且$n\le4096$追加三评分的串行/4/8流，共至多9唯一候选。
若没有新增的5至8直接用户亲和边且没有形成超过4成员的组，
则组划分、主要用户和同评分遍历与第27批相同，立即跳过重复整束搜索。
只保留亲和图构造和有界组检测的成本；没有将结构触发误称零成本。
排序专用边最多8倍原边，父节点消费对最多28个；128祖先上限、原图输出生命周期保持。
6个旧检查扩展为15项，包括8人组的环拒绝/拓扑保持以及小组不变时的重复跳过，全部通过。

979严格回归46.974秒、132独立标尺及21证书审计通过，CTest4/4 1.12秒，单文件编译无警告。
旧困难仍32/48、生命周期54/72，已观察36结构组由10至13成功；
恢复高扇出密集12阶段两个保留输出参数图和密集64阶段无保留输出一图，共同T/peak保持。
三次交替冻结27/28配对，各行含所有失败调用，时间单位ms：
\par\smallskip\noindent{\small\setlength{\tabcolsep}{3pt}\begin{tabular}{>{\raggedright\arraybackslash}p{3.2cm}>{\raggedleft\arraybackslash}p{1.5cm}>{\raggedleft\arraybackslash}p{3.1cm}>{\raggedleft\arraybackslash}p{3.1cm}>{\raggedleft\arraybackslash}p{1.8cm}}
集合 & 可行数 & 调用墙时ms & 线程CPU ms & CLI秒\\\hline
结构36 & 10/13 & 4378.474/\allowbreak{}4426.673 & 4203.125/\allowbreak{}4171.875 & 6.023/\allowbreak{}4.697\\
原困难开发24 & 17/17 & 5058.842/\allowbreak{}4764.959 & 4812.500/\allowbreak{}4437.500 & 5.801/\allowbreak{}5.249\\
生命周期第二种子36 & 27/27 & 2044.190/\allowbreak{}2336.549 & 2062.500/\allowbreak{}2312.500 & 3.634/\allowbreak{}3.090\\
\end{tabular}}\par\smallskip

结构共同10图调用墙时526.885/\allowbreak{}494.213、线程CPU531.250/\allowbreak{}484.375ms；
原困难共同17图1650.726/\allowbreak{}2029.004、1562.500/\allowbreak{}1984.375ms；生命周期共同27图1044.199/\allowbreak{}1243.233、1031.250/\allowbreak{}1171.875ms。
共同旧代码不触发新束却仍明显波动，线程CPU亦受本机实际频率、缓存竞争和计数粒度影响；
不能把全记录CPU下降作为确定加速，也不能隐去生命周期组全CPU增加12.1\%。
本批仅保留有限恢复及未解路径结构保护，没有普遍质量收益结论。

冻结源码\path{b52f3b8e2b54e0f8}及二进制后再生成种子20291028的12条未见高扇出图，
阶段数改为16/80/160，节点160至1659；局部/密集边、保留输出分别变化，
分类核池4/8/16/128、Stream4/8/16、Event0/3/$10^6$交叉配置，核请求按对应池上限收缩。
serial witness无Event、时长及peak不因核收缩改变，每图仍可证可行，输入未按solver成败过滤。
冻结27/28各运行一次均2/12通过，无新恢复，共同T/peak及模型保持；没有据未见结果继续调本批。
混合硬件、尺寸和种子的一次验证不能隔离单因素，但真实零增益说明宽组无法填补一般高扇出缺口。
单次求解墙时/线程CPU为2802.551/\allowbreak{}3099.255/\allowbreak{}2828.125/\allowbreak{}3109.375ms，仅诊断成本。
累计1147案例记录1080通过67个有witness仍未解；默认仍保持小消费束预算，后续应改善前沿定位与回溯。

\subsection{第二十九批：整组消费信用短名单}
大图目标组超过128时，原方法用组的一个任意代表节点估计所需输入和释放信用；
组中多个消费者共同完成后才能释放的输入可能全被漏算，已部分完成的组也有代表偏差。
新短名单遍历组中尚未完成成员，累计组输出、最大成员workspace，
按不同直接父节点去重汇总缺失分配；已完成父节点的剩余真实用户数等于组内剩余消费者数时，
把该输出计入预计释放，信用亦按不同父节点去重。
仍是排序近似而非内存可行性证明，真正束用原图逐节点分配/最后直接用户释放，再作完整Event模拟。

初稿分别尝试8人组和4人组的三评分，定位只有两个大图恢复；
默认撤回4人组的额外portfolio，并限制$512<n\le4096$、此前未解、实际组超过4人、初始目标组超过128，
最多9唯一候选。宽组是经验成本触发；初始目标组不超过128时，目标组数只能减少，
短名单永不截断，因此跳过是遍历不变的确定保护。
24项环拒绝、原拓扑、完整排列与alias检查覆盖513节点的新分支及小图不变短名单，通过。
979严格回归65.887秒、132最优标尺/21证书、CTest4/4 0.84秒，单文件编译无警告。

结构36由13至14成功，恢复1292节点高扇出密集保留输出图；混合硬件12由2至3成功，
恢复813节点密集无保留输出图；旧48困难32、新72生命周期54保持，共同T/peak保持。
三次交替冻结28/29，所有成功及失败调用成本如下：
\par\smallskip\noindent{\small\setlength{\tabcolsep}{3pt}\begin{tabular}{>{\raggedright\arraybackslash}p{3.2cm}>{\raggedleft\arraybackslash}p{1.5cm}>{\raggedleft\arraybackslash}p{3.1cm}>{\raggedleft\arraybackslash}p{3.1cm}>{\raggedleft\arraybackslash}p{1.8cm}}
集合 & 可行数 & 调用墙时ms & 线程CPU ms & CLI秒\\\hline
结构36 & 13/14 & 7765.223/\allowbreak{}4620.083 & 6953.125/\allowbreak{}4250.000 & 8.279/\allowbreak{}4.893\\
混合硬件12 & 2/3 & 3184.500/\allowbreak{}3752.623 & 2843.750/\allowbreak{}3312.500 & 3.366/\allowbreak{}3.941\\
原困难开发24 & 17/17 & 4163.003/\allowbreak{}3879.275 & 3828.125/\allowbreak{}3875.000 & 4.441/\allowbreak{}4.128\\
\end{tabular}}\par\smallskip

共同案例墙时/线程CPU为878.187/\allowbreak{}834.730/\allowbreak{}843.750/\allowbreak{}828.125、42.211/\allowbreak{}59.090/\allowbreak{}46.875/\allowbreak{}62.500、1768.485/\allowbreak{}1544.944/\allowbreak{}1531.250/\allowbreak{}1562.500（ms，前/后）。
混合硬件全线程CPU增加16.5\%，原困难增加1.2\%；结构组反而下降38.9\%，
但是只在旧搜索后追加候选没有减少前面算法步骤，该下降不支持计算量降低或确定加速。
结构组全CLI前有21.956/\allowbreak{}4.495/\allowbreak{}8.279秒，后有4.614/\allowbreak{}4.893/\allowbreak{}20.379秒，
共享宿主负载使配对绝对成本明显不稳；原始各轮数据保留，质量/可行性恢复可以复核，稳定时延收益尚未证明。

冻结源\path{1df00a5d1e229337}后另生成种子20301029的18条\emph{递归}共享图，
叶数16/64/256、节点46至766，平衡fork/join、随机不平衡切分、随机跨子树复用此前正输出生产者三种语法，
结合消费者输出0/小量保留，分类核请求1至8、时长1至10000，不依赖阶段窗口或生产者/消费者时长数量级分离。
递归生成顺序为独立拓扑witness，非局部边只引用已生成节点；显存限和deadline取其精确serial peak与总时长，
不按solver结果筛选。冻结28/29各一次均15/18通过，共同T/peak无变化，未用结果回调本批。
未解包括190节点非局部共享、766节点不平衡保留输出及766节点非局部保留输出三图。
全调用墙时/线程CPU为512.326/\allowbreak{}589.265/\allowbreak{}515.625/\allowbreak{}500.000ms，仅单次诊断，新增输入整体独立模拟/提交模型一致。
累计1165记录1097通过68个有witness仍未解；通过数含原解析不可行预期样本，不等于1097份可行计划。
记录间的共享官方样例和根输出参数控制亦不是独立随机拓扑，保留family分母而非用总比例掩盖困难分布。

\subsection{第三十批：超过72个sink的有界闭包投影}
较小宽消费图常有73至144个sink，原四人组投影直接跳过；新的回退仍只用于此前未解且$n\le512$。
按输出/扇出/执行时长的正输入亲和性合并sink，单组上限8，总组必须压至18以内，
随后沿用原直接依赖祖先闭包、完成覆盖mask及真实最后消费者计数。
相同完成闭包子集的live和直接用户数相同，失败mask可缓存；固定闭包内DFS顺序意味着并不完备枚举全部串行顺序。
维持800万基本探测预算，越界返回未找到，完整方案仍经独立模拟过滤。
只追加串行/4流/4流guard至多3候选，不把有界搜索失败当作证明无解。

新增三项解析宽闭包控制，sink数73/80/144，每8个消费者共享一个输出10的生产者，
生产者时长1，消费者时长100、workspace100、输出0，显存限110。
旧72上限拒绝，新投影所得serial $T=100N+\lceil N/8\rceil$、peak110逐项通过。
至少一个消费者启动时必须同时保留该父输出，peak110为必要界；两消费者不能并行，
打开两个未消费父输出会使消费者无法启动，故解析资源约束解释该图的窄前沿，不依赖样例ID。
三项属于单元解析控制，不额外混入benchmark记录数；已有24排序/alias检查继续通过。
979严格回归50.716秒、132最优标尺/21证书，CTest4/4 0.84秒、单文件编译无警告。

混合硬件12由3至5成功，恢复局部16阶段保留/不保留输出两图；结构36仍14、递归18仍15，旧困难32/48及生命周期54/72保持。
共同T/peak保持。三次交替冻结29/30，所有成功与失败调用成本列出：
\par\smallskip\noindent{\small\setlength{\tabcolsep}{3pt}\begin{tabular}{>{\raggedright\arraybackslash}p{3.2cm}>{\raggedleft\arraybackslash}p{1.5cm}>{\raggedleft\arraybackslash}p{3.1cm}>{\raggedleft\arraybackslash}p{3.1cm}>{\raggedleft\arraybackslash}p{1.8cm}}
集合 & 可行数 & 调用墙时ms & 线程CPU ms & CLI秒\\\hline
混合硬件12 & 3/5 & 3607.661/\allowbreak{}3615.847 & 3171.875/\allowbreak{}3171.875 & 4.318/\allowbreak{}3.808\\
结构36 & 14/14 & 4967.106/\allowbreak{}5274.692 & 4437.500/\allowbreak{}4328.125 & 5.664/\allowbreak{}6.326\\
递归18 & 15/15 & 725.652/\allowbreak{}552.109 & 562.500/\allowbreak{}515.625 & 1.101/\allowbreak{}0.900\\
\end{tabular}}\par\smallskip

共同线程CPU分别296.875/\allowbreak{}265.625、1218.750/\allowbreak{}1234.375、203.125/\allowbreak{}203.125ms；
未观察到稳定的普遍成本收益，也没有把计数量化下的相同值称作免费搜索。
全调用绝对时间继续受共享宿主影响，恢复案例与范围可复核，时延判断保留限制。
冻结源\path{6c9ea2d3b557c88a}后生成种子20311030共16图，阶段14/18、节点137至196，
局部/密集、输出保留0/1，以及原角色时长或所有算子独立1至100的两种时长分布；
各分布用不同种子，核池8/128、Stream4/16和Event0/3/$10^6$交叉设置，仍有独立serial witness。
冻结29/30各一次均13/16通过，无新增恢复，三个有解未解覆盖相近时长密集、保留输出密集及保留输出局部结构。
未按solver过滤、未据该结果回调30。全调用墙时768.291/\allowbreak{}685.341、线程CPU656.250/\allowbreak{}671.875ms，仅单次诊断。
累计1181记录1112按预期通过69个有witness仍未解，含原不可行预期样本及重复/参数记录的限制继续适用。

\subsection{第三十一批：严格半容量资源互斥下界}
面积界允许分数装箱，不能识别所有大请求必须串行的情况。
令$H_r=\{v:c_v>\lfloor C_r/2\rfloor\}$，类型$r$固定为一个AIC或AIV核池。
集合内任意两个请求之和大于该池容量，故其Compute区间不能重叠，
$T\ge\sum_{v\in H_r}d_v$。两个核池分别求界再取最大值，不将AIC与AIV集合的时长相加。

题面显存语义为\emph{从Compute开始同时分配本节点workspace与output}；
workspace在Compute结束释放，output至少存活至该Compute结束，之后按原直接消费者生命周期保留。
因此对$H_m=\{v:W_v+O_v>\lfloor M/2\rfloor\}$，整个Compute左闭右开区间内
均有$W_v+O_v$共同占用显存，同理$T\ge\sum_{v\in H_m}d_v$。
输出在Compute后的额外持留只会进一步限制调度，不将不同阶段、不共同存活的资源合成互斥关系。
独立公共模拟器逐时钟完成释放与分配，验证提交评估器的相同语义；分析界的证明不依赖solver模拟结果。
将三项与旧关键路径/面积/Stream下界取最大，只增强必要界及已有双界终止证书，不增加候选或搜索预算。

比较采用\texttt{request > capacity/2}避免乘2溢出；$W+O$、时长累加使用64位正数饱和加法。
旧核面积单项$d_vc_v$在已核对的$d_v\le10^9,c_v\le128$下至多$1.28\times10^{11}$，
累积继续饱和，不以32位值计算。严格大于一半不可换成大于等于：恰好一半允许两个任务重叠。
新增六项字面解析单元检查覆盖单池、分离双池、恰好半容量负控制、workspace互斥、
$W40+O20>M100/2$及三个$10^9$时长的64位累加；另保留288份旧分流递推、24项排序/环/alias和三份宽投影解析控制。
979严格回归65.071秒，无成败/T/同T峰值退步及模型差异；CTest4/4用时1.01秒，
仅\texttt{problem.h}+标准库的独立GCC编译无警告。

种子20321031构造12条解析benchmark控制：节点16/512/10000交叉AIC请求5/容量8、
AIV请求65/容量128、workspace$6\times10^8$/显存$10^9$及AIC恰好4/8负控制。
前三种的serial witness达到互斥时长和必需峰值界；恰好半容量的两个等长Stream达到核面积界。
稀疏ID、乱序记录、256可用Stream、$10^9$节点时长均使用实际JSON适配和独立事件模拟。
资源数值只用于标量模拟，未申请实际1GB设备显存。该族按理论专门构造，不能作为随机未见通用性能分布。
冻结30/31各三轮交替，12条均达到解析最优T/peak，证书从3份增至12份；
汇合旧120份四至六节点两流完整oracle及12份解析控制，144份独立最优标尺审计通过，33份命中双界证书。
小图oracle范围为本地显式直接同步模型，不代表更大图或所有合法冗余同步的全局穷举。

\par\smallskip\noindent{\small\setlength{\tabcolsep}{3pt}\begin{tabular}{>{\raggedright\arraybackslash}p{3.2cm}>{\raggedleft\arraybackslash}p{1.5cm}>{\raggedleft\arraybackslash}p{3.1cm}>{\raggedleft\arraybackslash}p{3.1cm}>{\raggedleft\arraybackslash}p{1.8cm}}
集合 & 可行数前/后 & 调用墙时ms & 线程CPU ms & CLI秒\\\hline
解析资源12 & 12/12 & 445.003/\allowbreak{}88.126 & 453.125/\allowbreak{}78.125 & 0.934/\allowbreak{}0.570\\
通用持出143 & 143/143 & 2812.660/\allowbreak{}3585.159 & 2687.500/\allowbreak{}2937.500 & 4.387/\allowbreak{}5.514\\
混合硬件12 & 5/5 & 3902.256/\allowbreak{}4119.915 & 3218.750/\allowbreak{}3250.000 & 4.555/\allowbreak{}4.412\\
\end{tabular}}\par\smallskip

表中成本包含\emph{全部失败搜索}，共同成功T/peak逐项不变，恢复数均0。
解析控制调用墙时下降80.2\%、线程CPU下降82.8\%；该效应来自达到证明界后跳过无收益搜索。
通用持出线程CPU增加9.3\%、调用墙时增加27.5\%，未观察泛化加速；
第一轮旧CLI耗时30.440秒，其余4.155/\allowbreak{}4.387秒，共享宿主抢占显著，不能挑该轮作为收益证据。
混合硬件全线程CPU增加1.0\%，共同5条线程CPU均406.250ms，负结果保留。
Win32线程CPU单项常按15.625ms量化，小记录0不表示实际零成本；三轮中位数也不是消除共享干扰的统计保证。
新界仅在线性扫描已有字段时累加，保留作为可证明低增量成本的专项终止优化，未据成本波动扩展搜索。

\subsection{最终冻结与覆盖缺口}
最终提交源SHA-256为\path{b68148bfed40aa1cd3e587efa4162c134e622b31ff2beac5425991de98e15b88}，二进制SHA-256为\path{5dfa5faef1f3a653cfeffabd7d0cd1ad6a9aec6d3617905976c944a1cd93a133}，
冻结到\path{research/batch31/}；第七批与所有早期批次源/数字仍独立保留。
在979条严格记录之后，另对202条困难记录执行一次最终功能核对，
33.814秒完整管线内133条返回合法计划、69条明确未找到；与冻结30的成败及共同T/peak相同。

\begin{center}\begin{tabular}{lrrr}
困难族 & 记录数 & 合法可行计划 & 有witness未解\\\hline
非连续消费者开发/原持出 & 48 & 32 & 16\\
生命周期开发/原持出 & 72 & 54 & 18\\
新结构更密边/高扇出 & 36 & 14 & 22\\
混合硬件与大Event & 12 & 5 & 7\\
递归共享图 & 18 & 15 & 3\\
sink与时长结构改变 & 16 & 13 & 3\\
合计 & 202 & 133 & 69\\
\end{tabular}\end{center}

最终1193条记录由严格979、困难202及新解析12组成，1124条\emph{按预期通过}。
其中一条压力输入被独立证明不满足deadline，预期异常处理通过；重复官方样例、根输出参数控制亦不能计为独立拓扑。
69条未解输入的真实可行witness已经独立验证，失败不是判定其无解，不伪报合法方案。
选择上述各记录的一次最终结果，调用墙时合计70.879秒、线程CPU64.094秒；
其中全部69条失败调用墙时19.334秒、线程CPU18.250秒。
这是一组\emph{非同期单次功能核对}，包含不同管线运行，不是新旧配对速度结论。
完整CLI分段为严格65.071秒、困难33.814秒、新解析首轮0.570秒，含JSON、witness验证与公共模拟的成本另于solver。
峰值宿主约405.5MiB来自通用143条JSON DOM及验证，不将它当设备显存或单文件solver专属内存。

所有成功计划经过公共独立模型与提交评估器核对；已知未解记录、错误、输入路径及各次墙时/线程CPU
见\path{reports/final-hard-batch31.json}。同模型核对和固定种子覆盖不证明未见随机DAG、任意核池/Stream数量、
超过已测10000节点、大图密边或所有Event竞争交叠的正确性/可行性覆盖。
当前没有官方包、隐藏得分、真实NPU执行或平台CPU/RAM时限，最终结论保持在本地题面模拟与明确oracle范围。

\subsection{后续工程工作量估算}
估算一名熟悉C++与调度算法的开发者，每工作日6小时有效研发，无额外NPU采购等待。
当前小图beam继续扩展到更一般交叠/随机长距离消费需要3至5日，含状态缓存、释放下界与新种子验证；
按实际阻塞补选择性Event需要2至4日，含防循环和同步删除实验；
搜索预算自适应与多轮局部改进需要2至3日；在已完成6节点oracle上扩更强状态/更大规模及时延对照需2至4日；
流式JSON与验证器单次指标返回需1至2日；取得官方包后的适配与验收需2至3日。
串行推进合计12至21工作日，约72至126有效小时；算法未知收益与官方接口差异会增加不确定性。
每个阶段先设独立witness可行率、共同T/peak及CPU预算，再决定是否保留，避免测试数量取代算法收益。

\begin{longtable}{>{\raggedright\arraybackslash}p{3.0cm}>{\raggedright\arraybackslash}p{5.2cm}>{\raggedright\arraybackslash}p{5.6cm}}
\toprule
方向 & 预期收益依据 & 下一阶段验收与边界\\\midrule
\endfirsthead
\toprule
方向 & 预期收益依据 & 下一阶段验收与边界\\\midrule
\endhead
\bottomrule
\endfoot
显存前沿beam/回溯 & 已有512节点以内宽度8搜索，任意长距离扇出仍可能没有可行候选 & 增强下界与状态预算，在新随机种子和更密DAG上比较恢复率、T、CPU，不宣称通用有解必达\\
选择性Event guard & 全局相邻guard可能添加多余同步，前沿guard保守预留会拖慢并行度 & 只对实测冲突补边，保持原始直接消费者释放；独立模拟排除循环/死锁及deadline回退\\
自适应搜索与局部改进 & 现有策略追加数量固定，初始邻域仅一次固定参考计划 & 学习有效策略预算，尝试多次改进、前后向压紧；用持出种子防止只拟合开发集\\
更强下界与精确oracle & 当前串行peak及两流6节点有标尺，更大并行图仍缺少最优对照 & 加入资源下界与状态缓存，给出更大规模最优差距；不可把可行witness当最优解\\
宿主测试成本 & JSON DOM约405MiB且公共模型接口每取一个指标重建模拟 & 流式逐case适配或一次取完整指标，分开测solver与测试管线收益，不改变官方提交接口\\
真实NPU落地 & 当前仅题面CPU模拟，没有硬件事件时延实测 & 获取官方测试包后接替本地schema；可选Ascend设备验证同步机制，保持题面模型与硬件测量分栏\\
\end{longtable}

当前有界启发式仍可能漏掉可行隐藏图；局部性回退限定至10000节点，Stream搜索上限256，
这些是本地实现取舍而非官方规模约束。代码不提供全局最优保证，真实设备行为也不能由这些CPU结果推出。
剩余工程余量主要在更强可行性搜索、选择性同步和下界/oracle；不宜为了增加测试数量而反复运行已通过的相同数据。

\section{复现入口与保存证据}
项目根目录运行\path{scripts/build.ps1}完成构建、CTest及样例。
\path{scripts/test.ps1}加\texttt{-Full}执行原压力集与独立提交编译。
研究数据由\path{scripts/generate_research_cases.py}生成。
\path{scripts/paired_benchmark.py}执行交替测量并保存输入/二进制/源代码SHA-256。
\path{scripts/research.ps1}串联两套研究测试、配对摘要和本报告生成。
\path{reports/paired-final-summary.json}与\path{reports/paired-holdout-final-summary.json}
为第六批数值证据入口，所有单次计划、CSV、失败输入保留。
冻结源码与二进制在\path{research/baseline-v1/}、\path{research/batch6/}与\path{research/batch7/}。
本报告为独立单文件，没有外部图片、bib或自定义字体依赖，使用XeLaTeX/ctexart编译。

\begin{thebibliography}{99}
\bibitem{official} Chaspark. 2026 Hackathon赛题27题面，核对日期2026-10-10.
\url{https://www.chaspark.com/#/races/competitions/1326715227709140992}.
\bibitem{heft} H. Topcuoglu, S. Hariri, M.-Y. Wu. Performance-effective and low-complexity task scheduling for heterogeneous computing.
IEEE TPDS, 13(3):260--274, 2002. DOI:10.1109/\allowbreak{}71.993206.
\url{https://disco.ethz.ch/courses/fs14/seminar/paper/Jochen/4.pdf}.
\bibitem{kolisch} R. Kolisch. Serial and parallel resource-constrained project scheduling methods revisited: Theory and computation.
EJOR, 90(2):320--333, 1996. \url{https://doi.org/10.1016/\allowbreak{}0377-2217(95)00357-6}.
\bibitem{hartmann} R. Kolisch, S. Hartmann. Experimental investigation of heuristics for resource-constrained project scheduling: An update.
EJOR, 174:23--37, 2006.
\url{https://www.hsba.de/fileadmin/user_upload/bereiche/_dokumente/6-forschung/profs-publikationen/Hartmann_2006_Experimental_investigation_of_Heuristics.pdf}.
\bibitem{memorydag} L. Marchal, H. Nagy, B. Simon, F. Vivien. Parallel scheduling of DAGs under memory constraints.
IPDPS, 2018.
\url{https://www.uni-bremen.de/fileadmin/user_upload/fachbereiche/fb3/infcon/CSLog/Publications/bsimon/papers/2018-MNSV-IPDPS.pdf}.
\bibitem{memorytree} L. Eyraud-Dubois, L. Marchal, O. Sinnen, F. Vivien. Parallel scheduling of task trees with limited memory.
arXiv:1410.0329, 2014. \url{https://arxiv.org/abs/1410.0329}.
\bibitem{cuda} NVIDIA. CUDA Driver API Reference Manual 13.4, Stream/Event Management，访问日期2026-10-11.
\url{https://docs.nvidia.com/cuda/cuda-driver-api/cuda_driver_api/group__CUDA__STREAM.html};
\url{https://docs.nvidia.com/cuda/cuda-driver-api/cuda_driver_api/group__CUDA__EVENT.html}.
\bibitem{ascend} Huawei Ascend. CANN 6.2.RC2，aclrtRecordEvent接口文档，访问日期2026-10-11.
\url{https://www.hiascend.com/doc_center/source/zh/canncommercial/62RC2/inferapplicationdev/aclcppdevg/aclcppdevg_03_0063.html}.
\bibitem{threadtime} Microsoft. GetThreadTimes function，Win32 API官方文档，访问日期2026-10-11.
\url{https://learn.microsoft.com/en-us/windows/win32/api/processthreadsapi/nf-processthreadsapi-getthreadtimes}.
\bibitem{fiedler} M. Fiedler. Algebraic connectivity of graphs. Czechoslovak Mathematical Journal, 23(2):298--305, 1973.
\url{https://dml.cz/handle/10338.dmlcz/101168}.
\end{thebibliography}
\end{document}
